\chapter{Supplementary Materials}
\label{app:supplementary}

\section{Preprocessing Pipeline}
\label{app:preprocessing}

The following steps were applied identically to both datasets:

\begin{enumerate}
    \item Load CSV data using \texttt{pandas.read\_csv()}.
    \item Remove duplicate records.
    \item Impute missing symptom values with 0 (symptom absent).
    \item Encode disease labels using \texttt{sklearn.preprocessing.LabelEncoder}.
    \item Split data into 80\% training / 20\% test using \texttt{train\_test\_split} with \texttt{stratify=y}.
    \item Apply \texttt{StandardScaler} to features for Logistic Regression and SVM only.
    \item Perform stratified 5-fold cross-validation for hyperparameter tuning.
\end{enumerate}

\section{Hyperparameter Search Spaces}
\label{app:hyperparameters}

Table~\ref{tab:hyperparameters} documents the hyperparameter search spaces used for grid search on Dataset~1 and the optimal values selected there. Dataset~2 uses fixed settings, which are listed after the table.

\begin{table}[H]
\centering
\caption{Hyperparameter search spaces and selected optimal values}
\label{tab:hyperparameters}
\small
\begin{tabular}{@{}llll@{}}
\toprule
\textbf{Classifier} & \textbf{Parameter} & \textbf{Search Space} & \textbf{Selected} \\
\midrule
Naive Bayes & alpha & \{0.01, 0.1, 0.5, 1.0\} & 0.01 \\
\midrule
\multirow{2}{*}{Logistic Regression} & C & \{0.01, 0.1, 1, 10, 100\} & 0.01 \\
& solver & \{lbfgs, saga\} & lbfgs \\
\midrule
\multirow{2}{*}{SVM (RBF)} & C & \{0.1, 1, 10, 100\} & 1 \\
& gamma & \{scale, auto, 0.01, 0.1\} & 0.01 \\
\midrule
\multirow{3}{*}{Random Forest} & n\_estimators & \{100, 200, 500\} & 100 \\
& max\_depth & \{10, 20, 50, None\} & 10 \\
& min\_samples\_split & \{2, 5, 10\} & 2 \\
\midrule
\multirow{4}{*}{XGBoost} & n\_estimators & \{100, 200, 500\} & 500 \\
& max\_depth & \{3, 6, 10\} & 3 \\
& learning\_rate & \{0.01, 0.1, 0.3\} & 0.3 \\
& subsample & \{0.8, 1.0\} & 0.8 \\
\bottomrule
\end{tabular}
\end{table}

For Dataset~2, the fixed settings are: Naive Bayes with alpha = 0.01, Logistic Regression with C = 0.01 and solver = lbfgs, SVM with hinge loss and alpha = 1e-4, Random Forest with 100 trees, max depth 10, and min_samples_split = 2, and XGBoost with 20 trees, max depth 3, learning rate 0.3, and subsample 0.8.

\section{Compute Environment}
\label{app:software}

All experiments were carried out on a single machine without GPU acceleration. The hardware and operating system configuration is summarized in Table~\ref{tab:hardware}, and the software stack in Table~\ref{tab:software}.

\begin{table}[H]
\centering
\caption{Hardware and operating system}
\label{tab:hardware}
\begin{tabular}{@{}ll@{}}
\toprule
\textbf{Component} & \textbf{Specification} \\
\midrule
Machine          & MacBook (2021 model) \\
Processor        & Apple M-series chip (ARM64) \\
Memory           & 16 GB unified memory \\
Storage          & SSD \\
GPU              & Not used (CPU-only training) \\
Operating System & macOS Sonoma \\
\bottomrule
\end{tabular}
\end{table}

\begin{table}[H]
\centering
\caption{Software environment}
\label{tab:software}
\begin{tabular}{@{}ll@{}}
\toprule
\textbf{Software} & \textbf{Version} \\
\midrule
Python     & 3.9+ \\
scikit-learn & 1.3+ \\
XGBoost    & 2.0+ \\
pandas     & 2.0+ \\
NumPy      & 1.24+ \\
matplotlib & 3.7+ \\
seaborn    & 0.12+ \\
joblib     & 1.3+ \\
jupyter    & 1.0+ \\
ipykernel  & 6.25+ \\
tqdm       & 4.66+ \\
\bottomrule
\end{tabular}
\end{table}

\section{Training Time Estimates}
\label{app:training_times}

Table~\ref{tab:training_times} shows the recorded wall-clock time (rounded) needed to fit the final model for each classifier on each dataset in the experiment runs. On Dataset~1, the overhead of the full grid search with 5-fold stratified cross-validation (Section~\ref{sec:hyperparameter_tuning}) is also added to the single-fit times, whereas on Dataset~2 the single fits with the fixed hyperparameters listed in Table~\ref{tab:hyperparameters} were evaluated instead, as exhaustive tuning was not computationally feasible at this scale. All the numbers vary with the system load.

\begin{table}[H]
\centering
\caption{Measured final-fit training times on the compute environment described in Table~\ref{tab:hardware}, rounded from the recorded experiment runs.}
\label{tab:training_times}
\begin{tabular}{@{}lll@{}}
\toprule
\textbf{Classifier} & \textbf{Dataset 1} (final fit) & \textbf{Dataset 2} (final fit) \\
\midrule
Naive Bayes         & $\sim$2 seconds  & $\sim$2 seconds \\
Logistic Regression & $\sim$5 seconds  & $\sim$3 minutes \\
SVM                 & $<$1 second (RBF) & $\sim$2.5 minutes (linear SGD) \\
Random Forest       & $\sim$9 seconds  & $\sim$3 seconds \\
XGBoost             & $\sim$25 seconds & $\sim$2.5 minutes \\
\midrule
\textbf{Total}      & $\sim$40 seconds & $\sim$8 minutes \\
\bottomrule
\end{tabular}
\end{table}

On Dataset~1, XGBoost (500 trees) and Random Forest dominate the fit times, whereas Naive Bayes and the small RBF SVM are so fast that they finish almost immediately. On Dataset~2 the picture is different: the linear models (Logistic Regression and the SGD-SVM) and XGBoost each need a few minutes to train, with the number of solver iterations or boosting rounds taking up most of the time over the 727-class output space; in comparison, Naive Bayes and the depth-capped Random Forest train in a few seconds.

